\section{سؤالی با قطعه‌کد (۱۵ نمره)}

قطعه‌کد زیر را بخوانید و بگویید خروجی‌اش چیست و چرا.

\begin{neocode}[neocodetitlefont=\small]{\lr{\pixellatin config.yaml}}
listen-on: 80

paths:
  "/":
    type: redirect
    target: "http://localhost:8080/index.html"
  "/dir/":
    type: servedir
    target: "/home/user/movies/"
\end{neocode}

\lr{\texttt{neocode}} متن نوار عنوان را به عنوان آرگومان اجباری می‌گیرد و هر
کلید \lr{\texttt{listings}} را در آرگومان اختیاری، مثلاً
\lr{\texttt{[listing options=\{language=Python\}]}}. داخل بلوک،
\lr{\texttt{@\ldots@}} به لاتک بازمی‌گردد، پس ماکروهایی مانند
\lr{\texttt{\textbackslash coursecode}} همچنان باز می‌شوند.

\subsection{فهرست تیک‌دار}

\begin{neochecklist}
    \citem موردی که انجام شده، با \lr{\texttt{\textbackslash citem}}
    \item موردی که انجام نشده، با \lr{\texttt{\textbackslash item}}
    \item موردی دیگر که پیش از تحویل باید انجام شود
\end{neochecklist}
